% =====================================================================
% frontmatter/folha-rosto.tex
% PT: Folha do júri (página 2 do modelo ESTSetúbal/UNIPS). Reproduz a
%     segunda página de TrabalhoUNIPS_TecnologiaSetubal_Mestrado_definitiva.docx.
%     NÃO precisa de editar este ficheiro: o texto vem de
%     config/metadata.tex. NÃO altere as posições \put nem as medidas em
%     bp. O sistema de coordenadas está explicado em frontmatter/cover.tex.
% EN: Jury page (page 2 of the ESTSetúbal/UNIPS model). Reproduces the
%     second page of TrabalhoUNIPS_TecnologiaSetubal_Mestrado_definitiva.docx.
%     You do NOT need to edit this file: the text comes from
%     config/metadata.tex. Do NOT change the \put positions or the bp
%     measurements. The coordinate system is explained in
%     frontmatter/cover.tex.
%
% PT: Medidas das caixas do Word (bp, a partir do canto superior esquerdo)
% EN: Word text-box measurements (bp, from the top-left corner)
%   autor / author   x = 85.05,  y = 263.4, largura/width 160.0, 16 pt, à direita/right-aligned
%   título / title   x = 249.95, y = 263.0, largura/width 290.1, 16 pt negrito/bold
%   linha / line     centro/centre x = 249.95, y = 263.5 -> 773.7, 3.75 pt, #CAF4FF
%   júri / jury      x = 249.95, y = 501.0, largura/width 275.1, 11 pt, entrelinha/leading 1.15
%   data / date      x = 250.15, y = 735.5, 10 pt
%
% PT: Nesta página o título pode ter até 12 linhas (a 13.ª encosta ao
%     bloco do júri, em y = 514.0); o limite prático é o da capa (10
%     linhas, ver frontmatter/cover.tex).
% EN: On this page the title can take up to 12 lines (a 13th line touches
%     the jury block, at y = 514.0); the practical limit is the cover's
%     (10 lines, see frontmatter/cover.tex).
% =====================================================================
\begin{titlepage}
\AddToHookNext{shipout/foreground}{%
  \setlength{\unitlength}{1bp}%
  \linespread{1}\fontfamily{phv}\fontseries{m}\fontshape{n}\selectfont
  \hyphenpenalty=10000 \exhyphenpenalty=10000
  % -------------------------------------------------------------------
  % PT: Autor: coluna da esquerda, 16 pt regular, maiúsculas, alinhado à direita.
  % EN: Author: left column, 16 pt regular, capitals, right-aligned.
  \put(92.25,-279.0){\parbox[t]{145.6bp}{\raggedleft\fontsize{16bp}{18.4bp}\selectfont
    \MakeUppercase{\authorname}}}%
  % -------------------------------------------------------------------
  % PT: Título (e subtítulo opcional): coluna da direita, 16 pt negrito, maiúsculas.
  % EN: Title (and optional subtitle): right column, 16 pt bold, capitals.
  \put(257.15,-278.6){\parbox[t]{275.7bp}{\raggedright\fontsize{16bp}{18.4bp}\selectfont\bfseries
    \MakeUppercase{\thesistitle}\ifdefempty{\thesissubtitle}{}{\MakeUppercase{: \thesissubtitle}}}}%
  % -------------------------------------------------------------------
  % PT: Linha azul vertical #CAF4FF (no Word: linha de 510.2 pt rodada 270 graus).
  % EN: Vertical blue line #CAF4FF (in Word: 510.2 pt line rotated 270 degrees).
  \put(248.075,-773.7){\textcolor[HTML]{CAF4FF}{\rule{3.75bp}{510.2bp}}}%
  % -------------------------------------------------------------------
  % PT: Júri, 11 pt, entrelinha 1,15. Uma linha por membro: Presidente,
  %     Orientador e um ou dois vogais. A linha do segundo vogal só aparece
  %     se \jurysecondvogal não estiver vazio (ver config/metadata.tex).
  % EN: Jury, 11 pt, 1.15 line spacing. One line per member: president,
  %     supervisor and one or two examiners ("vogais"). The second
  %     examiner's line only appears if \jurysecondvogal is not empty
  %     (see config/metadata.tex).
  \put(257.15,-514.0){\parbox[t]{260.7bp}{\raggedright\fontsize{11bp}{14.55bp}\selectfont
    \setlength{\parindent}{0pt}\setlength{\parskip}{0pt}%
    {\bfseries JÚRI\par}
    \mbox{}\par
    Presidente: \jurypresident\par
    \juryadvisorlabel: \juryadvisor\par
    Vogal: \juryvogal\par
    \ifdefempty{\jurysecondvogal}{}{Vogal: \jurysecondvogal\par}}}%
  % -------------------------------------------------------------------
  % PT: Data da realização da prova (mês e ano), 10 pt.
  % EN: Date of the public defence (month and year), 10 pt.
  \put(257.35,-747.4){\parbox[t]{230.5bp}{\raggedright\fontsize{10bp}{13.2bp}\selectfont
    \defencemonth\ de \defenceyear}}%
}
% PT: Página vazia: o conteúdo é todo desenhado pelo bloco acima.
% EN: Empty page: all content is drawn by the block above.
\null
\end{titlepage}
